\documentclass{tutodoc}

\usepackage{../preamble.cfg}


\begin{document}

\section{Reuse from...}

Here are the key points to remember to use palettes offered by projects listed in the section \ref{projects-used}.
%
\begin{enumerate}
    \item Shifted palettes%
    \footnote{
        By design, \colorcet\ proposes several shifted palettes.
}
    and reversed ones%
    \footnote{
        \matplotlib\ palettes with names ending in \tdoccodein{py}+_r+ are reversed versions.
}
    are not included in \thisproj.
    In rare cases, some palettes are excluded because they appear redundant or are too large.
    All excluded palettes are listed in Appendix \ref{ignored-palettes} (page \pageref{ignored-palettes}).


	\item
	\thisproj\ uses 'semantic' \verb+CamelCase+ naming. Therefore, palette names such as
	\tdoccodein{text}{berlin}, \tdoccodein{text}{gist_heat} and \tdoccodein{text}{Pinkgrey}
	become
	\tdoccodein{text}{Berlin}, \tdoccodein{text}{GistHeat} and \tdoccodein{text}{PinkGrey}.
	To avoid naming collisions, certain palettes have been renamed : a full list of these changes is provided in Appendix \ref{renamed-palettes} (page \pageref{renamed-palettes}).
	You can see the full list of names in Appendix \ref{all-palette-names} (page \pageref{all-palette-names}).
\end{enumerate}


% -------------------- %


\begin{tdoccaut}
    Most \thisproj\ implementations add the \verb+pal+ prefix to standardized \verb+CamelCase+ names.
    See the section \ref{products-all}.
\end{tdoccaut}


\begin{tdocnote}
    Most \thisproj\ implementations provide methods to easily obtain reversed palettes, sub-palettes, and color-shifted palettes.
    See the section \ref{products-all}.
\end{tdocnote}

\end{document}
